% ==========================================
% BAB II STUDI LITERATUR
% ==========================================
\chapter{STUDI LITERATUR}
\label{chap:studi-literatur}

\section{Enterprise Resource Planning (ERP)}

\subsection{Definisi dan Evolusi ERP}

Enterprise resource planning (ERP) adalah sistem informasi yang didesain untuk integrasi dan optimisasi proses bisnis dan transaksi di sebuah enterprise dengan mencakup area berbeda seperti operasional (entah itu manufaktur atau hanya sekadar inventori), sumber daya manusia (SDM), keuangan/akutansi, penjualan, dan lainnya \cite{boza2015}. ERP sudah terbukti secara efektif meningkatkan efisiensi operasional dan membantu pengambilan keputusan di perusahaan kecil maupun besar \cite{hasanah2024}. Di "The Evolution of ERP Systems: A Literature Review" Justin Goldston membagi keuntungan penggunaan ERP menjadi tiga kategori keuntungan teknologi, keuntungan knowledge sharing, dan keuntungan kepemimpinan. Keuntungan teknologi mencakup keunggulan kompetitif, peningkatan performa finansial, Kemitraan Rantai Pasokan, dan pemahaman pelanggan. Keuntungan knowledge sharing mencakup kemudahan pertukaran informasi antar unit organisasi, Meningkatkan pembelajaran organisasi selama implementasi, dan memungkinkan departemen untuk berkolaborasi alih-alih menggunakan sistem berbeda. Keuntungan kepemimpinan mencakup Memungkinkan keterlibatan dan dukungan manajemen, Meningkatkan komunikasi antar kelompok bisnis, Memungkinkan penyelesaian konflik antar departemen, dan mendukung praktik improvement berkelanjutan \cite{goldston2020evolution}.

Sebelum ERP menjadi sistem informasi yang kita kenal sekarang, sistem informasi ini mengalami evolusi yang awalnya dari aplikasi kontrol inventori hingga menjadi alat bantu mengambil keputusan dan automasi proses bisnis. Berikut sejarah singkat evolusi ERP:

Pada awal era komputer (1960-an), banyak organisasi yang mulai memperkenalkan komputer dalam operasional mereka dan mengembangkan aplikasi sederhana untuk melacak inventaris, membantu dalam pemesanan material, dan mengatur produksi barang jadi. Pendekatan ini, yang dikenal sebagai "kontrol inventori", merupakan langkah awal perusahaan dalam menjalankan sisi operasional organisasi secara sistematis dan terkomputerisasi. Dari kontrol inventori ini mulailah muncul ide \textit{Materials Requirements Planning} (MRP), versi yang lebih terintegrasi dari "kontrol inventori" saja. Pada tahun 1970-an aplikasi MRP mulai dikenalkan dengan kemampuan yang memungkinkan organisasi melakukan pembelian, peramalan, dan penjadwalan produksi. Pada dekade ini lahir perusahaan vendor ERP di antaranya SAP dan J.D. Edwards \cite{goldston2020evolution}.

Pada 1980-an J.D. Edwards meningkatkan aplikasi MRP dengan menambahkan fitur, di antaranya \textit{closed-loop scheduling}, \textit{enhanced shop floor reporting}, dan \textit{forward scheduling}. Peningkatan ini dipromosikan dengan nama MRP-II. Pada dekade ini vendor bertambah dalam jumlah dan ukuran, di antaranya SAP, IBM, J.D. Edwards, Baan, PeopleSoft, dan Oracle. Dengan kemajuan MRP-II, organisasi mulai menggunakan teknologi untuk membantu mengambil keputusan. Pada dekade 1990-an integrasi MRP diperluas hingga ke operasional organisasi dan akuntansi perusahaan. Gartner Group secara resmi menciptakan istilah \textit{Enterprise Resource Planning} (ERP). Dengan ketidakpastian Y2K, perusahaan berlomba-lomba untuk menginstal aplikasi ERP menyebabkan ledakan pada ERP. Ledakan ini diikuti dengan perkembangan keenam vendor utama \cite{goldston2020evolution}.

Pada 2001 gelembung \textit{dotcom} (\textit{dotcom bubble}) terjadi, yang menyebabkan saham dan kapitalisasi pasar perusahaan vendor mengalami penurunan drastis hingga terjadi banyak akuisisi. Oracle mengakuisisi J.D. Edwards dan PeopleSoft sedangkan Infor Global Solutions mengakuisisi Baan dan produk MAPICS IBM, sehingga menyisakan 3 vendor raksasa Oracle, SAP dan Infor yang sampai sekarang masih menjadi vendor utama perusahaan besar. Setelah beberapa tahun, dikenalkan \textit{cloud computing} yang memindahkan tanggung jawab perangkat keras ke vendor, mengurangi overhead IT. Pada tahun 2016, laporan ERP menunjukkan peningkatan sebesar 40 persen dibandingkan tahun 2015 pada perusahaan yang mengimplementasikan \textit{cloud}. Sekarang pengembangan ERP berfokus pada pengurangan pemborosan dalam operasi dan terminologi baru diciptakan yaitu \textit{Sustainable Enterprise Resource Planning} (S-ERP) \cite{goldston2020evolution}.

\subsection{Komponen dan Arsitektur ERP}

Sistem ERP memiliki arsitektur berlapis (\textit{multi-tier architecture}) yang terdiri dari lapisan presentasi (\textit{presentation layer}), lapisan aplikasi (\textit{application layer}), dan lapisan basis data (\textit{database layer}). Lapisan presentasi menangani antarmuka pengguna dan interaksi dengan sistem, lapisan aplikasi memproses logika bisnis dan aturan enterprise, sedangkan lapisan basis data mengelola penyimpanan dan pengambilan data. Berbeda dengan aplikasi web sederhana, ERP memiliki kompleksitas teknis yang tinggi karena harus mengintegrasikan berbagai fungsi bisnis yang berbeda dalam satu sistem terpadu dengan menjaga konsistensi dan integritas data di seluruh organisasi \cite{boza2015}.

Komponen utama ERP terdiri dari berbagai modul fungsional yang mencakup modul keuangan dan akuntansi, modul manajemen inventori, modul sumber daya manusia (\textit{human resources management}), modul manajemen hubungan pelanggan (\textit{customer relationship management}/CRM), modul rantai pasokan (\textit{supply chain management}/SCM), dan modul manufaktur. Modul-modul ini saling terintegrasi melalui basis data terpusat dan mekanisme pertukaran data yang memungkinkan informasi mengalir secara otomatis antar fungsi bisnis. Integrasi ini merupakan salah satu nilai utama ERP dibandingkan dengan penggunaan sistem-sistem terpisah yang tidak terkoneksi, karena memungkinkan organisasi untuk memiliki satu sumber kebenaran (\textit{single source of truth}) dan mengurangi redundansi data \cite{boza2015}\cite{goldston2020evolution}.

Sebagai tulang punggung sistem ERP, infrastruktur basis data harus memenuhi persyaratan performa yang tinggi, menjaga konsistensi data, dan mendukung transaksi konkuren dari berbagai modul secara bersamaan. Basis data ERP harus mampu menangani volume data yang besar dengan waktu respons yang cepat, menjamin integritas referensial antar tabel, serta menyediakan mekanisme \textit{backup} dan \textit{recovery} untuk keandalan sistem. Dalam perkembangan modern, arsitektur ERP juga mengadopsi pendekatan terdistribusi dan berorientasi layanan (\textit{service-oriented architecture}/SOA) yang memungkinkan fleksibilitas deployment dan integrasi dengan sistem eksternal melalui antarmuka pemrograman aplikasi (API).

\subsection{Contoh-Contoh ERP}

Seperti yang telah dijelaskan sebelumnya, \textit{vendor} ERP komersial seperti Oracle, SAP, dan Infor mendominasi pasar \textit{enterprise} besar dengan sistem yang matang dan kaya fitur \cite{goldston2020evolution}. Solusi dari \textit{vendor} besar ini menawarkan stabilitas produk yang tinggi, dukungan teknis yang komprehensif, dan dokumentasi yang ekstensif. Namun, keuntungan tersebut datang dengan biaya implementasi yang sangat tinggi, lisensi mahal berbasis \textit{per user} atau transaksi, kompleksitas implementasi yang membutuhkan tim ahli khusus, dan risiko \textit{vendor lock-in}. Faktor-faktor ini membuat solusi komersial kurang praktis bagi organisasi dengan anggaran terbatas.

Sebagai alternatif yang lebih terjangkau, sistem ERP \textit{open-source} muncul dengan keuntungan yang menarik. Tidak ada biaya lisensi, \textit{source code} tersedia untuk modifikasi, komunitas pengguna yang aktif menyediakan dukungan, dan arsitektur yang modular memungkinkan organisasi mengimplementasikan hanya modul yang relevan dengan kebutuhan bisnis mereka. Pendekatan modular ini penting untuk organisasi yang tidak memerlukan fungsionalitas \textit{enterprise} penuh tetapi membutuhkan sistem yang dapat berkembang seiring pertumbuhan bisnis \cite{boza2015}. Dengan menghilangkan hambatan biaya dan meningkatkan fleksibilitas, solusi \textit{open-source} membuka akses ERP bagi organisasi yang sebelumnya tidak memiliki sumber daya untuk digitalisasi.

Di antara berbagai ERP \textit{open-source} yang tersedia, Odoo dan ERPNext adalah yang paling matang dan banyak diadopsi. Odoo merupakan ERP \textit{open-source} paling populer dengan komunitas global yang besar, ekosistem aplikasi yang luas melalui \textit{Odoo App Store}, dan dokumentasi yang detail \cite{odoo2025documentation}. Arsitektur modular Odoo memungkinkan organisasi memilih dan mengintegrasikan modul sesuai kebutuhan, menjadikannya fleksibel untuk berbagai skala bisnis. ERPNext dikembangkan di atas \textit{framework} Frappe menggunakan Python, dirancang spesifik untuk usaha kecil dan menengah dengan antarmuka intuitif dan implementasi yang lebih ringan \cite{erpnext2024introduction}. ERPNext mengedepankan kesederhanaan tanpa mengorbankan fungsionalitas inti dan telah mendapat adopsi yang meningkat termasuk di Indonesia. Selain keduanya, terdapat alternatif lain seperti Apache OFBiz dan Dolibarr. Odoo dan ERPNext relevan untuk penelitian ini karena kombinasi antara maturitas, dokumentasi yang baik, arsitektur modular, dan kemampuan \textit{deployment} pada infrastruktur beragam—menjadikan keduanya kandidat ideal untuk eksplorasi pada perangkat hemat daya seperti \textit{single-board computer}.

\section{Model Deployment Infrastruktur}

Pemilihan model \textit{deployment} infrastruktur merupakan keputusan strategis yang berdampak langsung pada biaya operasional, kontrol sistem, keamanan data, dan skalabilitas aplikasi \textit{enterprise} seperti ERP. Model \textit{deployment} menentukan bagaimana infrastruktur komputasi, penyimpanan, dan jaringan diatur, dikelola, dan diakses oleh organisasi \cite{Younus2024}. Dalam konteks ERP untuk UMKM, pemilihan model yang tepat sangat krusial karena organisasi skala kecil dan menengah memiliki keterbatasan anggaran, sumber daya teknis, dan toleransi risiko yang berbeda dari \textit{enterprise} besar. Secara umum, terdapat tiga model utama yang mendominasi: \textit{on-premise}, \textit{cloud}, dan \textit{hybrid}—masing-masing dengan karakteristik, keuntungan, dan tantangan yang perlu dipahami untuk menentukan kelayakan infrastruktur alternatif seperti perangkat hemat daya \cite{odunayo2018cloud}.

\subsection{On-Premise}

Model \textit{on-premise} merujuk pada \textit{deployment} infrastruktur secara internal di lokasi fisik yang dimiliki atau disewa oleh organisasi, di mana seluruh perangkat keras (\textit{hardware}), perangkat lunak (\textit{software}), dan manajemen sistem menjadi tanggung jawab penuh organisasi tersebut \cite{Younus2024}. Dalam model ini, organisasi harus menginvestasikan modal awal yang besar untuk membeli \textit{server}, perangkat jaringan, sistem penyimpanan, lisensi perangkat lunak, serta membangun atau menyewa ruang \textit{data center} dengan infrastruktur pendukung seperti pendingin udara dan \textit{uninterruptible power supply} (UPS). Organisasi juga bertanggung jawab penuh atas instalasi, konfigurasi, pemeliharaan, pembaruan sistem, pencadangan data (\textit{backup}), dan perencanaan pemulihan bencana (\textit{disaster recovery}).

Keuntungan utama model ini terletak pada kontrol penuh terhadap data, konfigurasi sistem, dan jadwal pembaruan yang memungkinkan organisasi menyesuaikan infrastruktur dengan kebutuhan spesifik tanpa terikat pada batasan penyedia layanan eksternal. Keamanan fisik dapat dijaga langsung, dan \textit{uptime} sistem tidak bergantung pada koneksi internet. Dari perspektif \textit{Total Cost of Ownership} (TCO) jangka panjang, model ini dapat lebih ekonomis karena biaya lisensi umumnya dibayar sekali atau dengan skema pembaruan tahunan yang lebih murah dibandingkan biaya langganan \textit{cloud} yang terus-menerus \cite{Younus2024}. Namun, investasi awal (\textit{capital expenditure}/CAPEX) yang besar menciptakan hambatan finansial tinggi, terutama untuk organisasi dengan anggaran terbatas. Tanggung jawab penuh atas pemeliharaan sistem memerlukan sumber daya manusia dengan keahlian teknis memadai, dan skalabilitas menjadi tantangan karena peningkatan kapasitas memerlukan pembelian perangkat keras baru yang tidak dapat dilakukan secara instan.

\subsection{Cloud}

Konsep \textit{cloud computing} mengubah paradigma \textit{deployment} infrastruktur dengan menawarkan akses \textit{on-demand} ke \textit{shared pool} sumber daya komputasi melalui internet, menghilangkan kebutuhan organisasi untuk memiliki dan mengelola perangkat keras secara langsung \cite{odunayo2018cloud}. Model ini dibangun di atas lima karakteristik fundamental: \textit{on-demand self-service}, \textit{broad network access}, \textit{resource pooling}, \textit{rapid elasticity}, dan \textit{measured service} yang menyediakan transparansi penggunaan untuk penagihan berbasis konsumsi aktual \cite{odunayo2018cloud}\cite{Younus2024}. Sejak diperkenalkan oleh Salesforce pada 1999 dengan model SaaS, diikuti Amazon Web Services (AWS) pada 2006 dengan layanan IaaS, \textit{cloud computing} telah berkembang menjadi ekosistem kompetitif yang melahirkan tiga model layanan utama: \textit{Infrastructure-as-a-Service} (IaaS), \textit{Platform-as-a-Service} (PaaS), dan \textit{Software-as-a-Service} (SaaS).

    \subsubsection{IaaS (Infrastructure-as-a-Service)}

Model IaaS menyediakan infrastruktur komputasi fundamental berupa \textit{server virtual}, penyimpanan (\textit{storage}), jaringan, dan sistem operasi sebagai layanan yang dapat disewa sesuai kebutuhan \cite{odunayo2018cloud}. Penyedia layanan \textit{cloud} (\textit{Cloud Service Provider}/CSP) bertanggung jawab mengelola \textit{data center} fisik, perangkat keras, virtualisasi, dan jaringan, sementara pengguna memiliki kontrol penuh atas sistem operasi, \textit{middleware}, dan aplikasi \cite{Younus2024}. Karakteristik utama IaaS mencakup skalabilitas tinggi yang dapat disesuaikan secara otomatis (\textit{auto-scaling}), model pembayaran \textit{pay-as-you-go} yang mengubah CAPEX menjadi \textit{operational expenditure} (OPEX), dan dukungan untuk \textit{Infrastructure as Code} (IaC) yang memungkinkan otomasi dan replikasi infrastruktur secara konsisten. Contoh layanan populer meliputi Amazon EC2, Microsoft Azure Virtual Machines, Google Cloud Compute Engine, dan layanan penyimpanan seperti Amazon S3.

    \subsubsection{PaaS (Platform-as-a-Service)}

PaaS menyediakan platform pengembangan aplikasi yang mencakup sistem operasi, \textit{runtime environment}, basis data, dan \textit{development tools} sehingga pengembang dapat fokus pada logika aplikasi tanpa mengelola infrastruktur yang mendasarinya \cite{odunayo2018cloud}. Penyedia layanan mengelola seluruh \textit{stack} infrastruktur hingga lapisan platform, sementara pengguna hanya bertanggung jawab atas kode aplikasi dan data \cite{Younus2024}. Karakteristik PaaS mencakup \textit{built-in tools} untuk pengembangan, pengujian, dan \textit{deployment}, mekanisme keamanan terintegrasi, kemampuan integrasi melalui API, serta skalabilitas otomatis yang disesuaikan dengan beban aplikasi. Contoh layanan yang banyak digunakan adalah AWS Elastic Beanstalk, Microsoft Azure App Service, dan Google App Engine yang memungkinkan tim pengembangan mempercepat waktu peluncuran produk (\textit{time-to-market}) dengan mengurangi overhead operasional.

    \subsubsection{SaaS (Software-as-a-Service)}

SaaS menyediakan aplikasi perangkat lunak yang di-\textit{hosting} oleh penyedia layanan dan diakses melalui internet tanpa perlu membeli lisensi, menginstal, atau memelihara aplikasi secara lokal \cite{odunayo2018cloud}. Penyedia layanan bertanggung jawab penuh atas seluruh \textit{stack} teknologi mulai dari infrastruktur fisik hingga aplikasi itu sendiri, termasuk pembaruan, keamanan, dan ketersediaan layanan \cite{Younus2024}. Karakteristik utama SaaS meliputi arsitektur \textit{multi-tenancy} yang melayani banyak pengguna dengan data terisolasi secara logis, model \textit{on-demand}, skalabilitas otomatis, serta kemampuan akses dari mana saja selama terhubung internet. Contoh aplikasi SaaS mencakup Google Workspace, Microsoft 365, Salesforce, serta Odoo \textit{Online} dan ERPNext \textit{Cloud} dalam konteks ERP yang menyediakan kemudahan implementasi dengan waktu \textit{deployment} minimal, namun dengan keterbatasan dalam kustomisasi mendalam dibandingkan \textit{deployment} \textit{on-premise} atau IaaS.

\subsection{Hybrid}

Model \textit{hybrid cloud} menggabungkan infrastruktur \textit{private cloud} (yang dapat berupa \textit{on-premise} atau \textit{hosted private cloud}) dengan \textit{public cloud} melalui teknologi standar yang memungkinkan portabilitas data dan aplikasi antar lingkungan \cite{odunayo2018cloud}. Pendekatan ini memungkinkan organisasi mengalokasikan beban kerja (\textit{workload}) secara strategis: aplikasi inti yang sensitif, data yang tunduk pada regulasi ketat, atau sistem yang memerlukan kontrol penuh dapat dijalankan di infrastruktur \textit{private}, sementara beban kerja yang memerlukan skalabilitas dinamis atau aplikasi dengan fluktuasi permintaan tinggi dapat di-\textit{deploy} di \textit{public cloud} \cite{Younus2024}.

Keuntungan model \textit{hybrid} mencakup fleksibilitas dalam alokasi sumber daya yang memungkinkan organisasi menyesuaikan kapasitas secara dinamis tanpa investasi berlebih, optimisasi biaya dengan menjalankan beban kerja stabil di infrastruktur \textit{private} sambil memanfaatkan \textit{public cloud} untuk kebutuhan sementara, serta kemampuan menjaga kontrol dan keamanan tinggi untuk data sensitif sambil tetap menikmati skalabilitas layanan \textit{public cloud}. Model ini juga mendukung kepatuhan terhadap regulasi (\textit{compliance}) yang mewajibkan data tertentu disimpan di lokasi spesifik \cite{odunayo2018cloud}\cite{Younus2024}. Namun, kompleksitas model \textit{hybrid} membawa tantangan tersendiri: manajemen infrastruktur tersebar memerlukan keahlian teknis tinggi dan \textit{tools} orkestrasi canggih, integrasi antar lingkungan dapat menimbulkan masalah latensi dan risiko keamanan pada \textit{data transfer}, serta konsistensi konfigurasi dan kebijakan tata kelola memerlukan standarisasi dan otomasi untuk menghindari \textit{downtime} atau celah keamanan. Meskipun demikian, prediksi Gartner pada tahun 2017 menunjukkan bahwa hingga 50 persen perusahaan \textit{mainstream} akan mengadopsi model \textit{hybrid cloud} karena keseimbangan yang ditawarkan antara kontrol, fleksibilitas, dan efisiensi biaya \cite{odunayo2018cloud}.

\section{Arsitektur Prosesor dan Single-Board Computer (SBC)}

\subsection{Arsitektur Prosesor CISC/x86 vs RISC/ARM}

Arsitektur prosesor merupakan fondasi utama bagi kinerja dan efisiensi perangkat komputasi, mulai dari server kelas enterprise hingga perangkat hemat daya seperti \textit{single-board computer} (SBC). Dua paradigma utama dalam desain arsitektur prosesor adalah \textit{Complex Instruction Set Computer} (CISC) yang umum diimplementasikan pada x86, dan \textit{Reduced Instruction Set Computer} (RISC) yang identik dengan ARM. Setiap paradigma menawarkan pendekatan berbeda dalam pengelolaan instruksi serta sumber daya perangkat keras~\cite{blem2013power}.

Arsitektur CISC/x86 mengutamakan kompleksitas instruksi sehingga satu instruksi mampu menginisiasi rangkaian operasi sekaligus. Instruksi pada x86 bersifat variabel, memiliki banyak mode pengalamatan, serta penanganan operasi memori yang fleksibel. Komputer berbasis x86 biasa digunakan pada desktop dan server karena mendukung optimalisasi performa melalui \textit{out-of-order execution} serta teknik mikroarsitektur lainnya~\cite{aroca2012towards}. Contoh implementasinya antara lain Intel Core i7 dan AMD Ryzen.

Sebaliknya, arsitektur RISC/ARM menitikberatkan desain instruksi yang sederhana dan seragam dengan panjang tetap sehingga eksekusi instruksi dapat lebih efisien. Prosesor ARM umum ditemukan pada perangkat hemat daya seperti smartphone, tablet, dan SBC. Beberapa tahun terakhir, ARM berkembang pesat dan mulai memasuki pasar server berkat efisiensi konsumsi energi serta biaya total kepemilikan yang lebih rendah~\cite{aroca2012towards}. Contoh prosesor ARM populer meliputi Cortex-A9, Cortex-A76, dan ARMv8.

Walaupun mendekati kontras secara konseptual, optimalisasi mikroarsitektur modern telah memperkecil perbedaan performa maupun efisiensi antara CISC/x86 dan RISC/ARM. Studi mutakhir menunjukkan efisiensi energi serta performa lebih bergantung pada teknik desain prosesor serta optimasi perangkat keras dan perangkat lunak, bukan sekadar tipe arsitektur~\cite{blem2013power}. Oleh karena itu, preferensi arsitektur kini lebih didasarkan pada kebutuhan aplikasi serta ekosistem perangkat.

\begin{table}[ht]
\centering
\begin{tabular}{|l|c|c|}
\hline
\textbf{Karakteristik} & \textbf{CISC/x86} & \textbf{RISC/ARM} \\
\hline
Desain Instruksi & Variabel, kompleks & Tetap, sederhana \\
Eksekusi & Multi-siklus/out-of-order & Single-siklus/in-order \\
Konsumsi Daya & Cenderung tinggi & Lebih rendah \\
Aplikasi Utama & PC, Server & Smartphone, Tablet, SBC \\
Upgrade Performa & Optimalisasi mikroarsitektur & Efisiensi dan integrasi SoC \\
\hline
\end{tabular}
\caption{Perbandingan Arsitektur CISC/x86 dan RISC/ARM}
\label{tab:risc-cisc}
\end{table}


Pada praktiknya, CISC/x86 dipilih untuk skenario yang membutuhkan kompatibilitas software legacy dan performa tinggi, seperti data center atau workstation. RISC/ARM lebih unggul dalam aplikasi hemat daya, edge computing, serta lingkungan dengan tuntutan efisiensi biaya total kepemilikan~\cite{aroca2012towards}. Fleksibilitas arsitektur memungkinkan vendor mengoptimalkan implementasi sesuai kebutuhan pasar.

\subsection{Single-Board Computer (SBC): Raspberry Pi dan Orange Pi}

\textit{Single-board computer} (SBC) adalah komputer utuh yang seluruh komponen utamanya—prosesor, memori, penyimpanan, dan port I/O—terintegrasi pada satu papan sirkuit. Konsep SBC mulai dikenal luas sejak era 1970-an dengan munculnya produk seperti KIM-1 dan Apple I, yang menawarkan solusi komputasi sederhana dan terjangkau untuk keperluan pendidikan dan hobi elektronika. Namun, perkembangan pesat SBC terjadi pada dekade 2010-an berkat inovasi teknologi dan kebutuhan terhadap perangkat hemat daya serta ukuran ringkas yang mendukung aplikasi embedded dan edge computing.

Seiring waktu, SBC berevolusi dari sekadar perangkat pembelajaran menuju platform fleksibel—dari robotika hingga digital signage dan otomasi industri. Faktor harga ekonomis, komunitas pengguna yang berkembang, serta support ekosistem software membuat SBC semakin diminati, terutama oleh pengembang yang membutuhkan prototyping cepat dan skalabilitas rendah. Dua contoh SBC yang paling sering digunakan adalah Raspberry Pi dan Orange Pi.

Raspberry Pi 5 merupakan generasi terbaru dari keluarga Raspberry Pi yang menawarkan peningkatan substansial dalam performa, konektivitas, dan dukungan multimedia. Produk ini dilengkapi prosesor quad-core ARM Cortex-A76 2.4GHz, RAM hingga 8GB, dual HDMI 4K, PCIe standar, serta berbagai port baru yang mendukung aplikasi AI, otomasi, dan multimedia terbaru. Keunggulan Raspberry Pi 5 tidak hanya pada performa yang jauh lebih baik dari generasi sebelumnya, namun juga pada kestabilan software dan dokumentasi yang sangat lengkap~\cite{raspberrypi2024pi5}.

Orange Pi 5 Pro adalah SBC rilisan terbaru dari Xunlong yang membidik segmen high-performance embedded computing dengan harga terjangkau. Menggunakan prosesor octa-core Rockchip RK3588S hingga 2.4GHz dan RAM LPDDR4/LPDDR5 hingga 32GB, Orange Pi 5 Pro mendukung output video 8K, slot M.2 NVMe SSD, serta berbagai fitur konektivitas lanjutan. Keunggulan utama Orange Pi 5 Pro terletak pada fleksibilitas hardware serta komunitas yang mulai berkembang, sehingga menjadi alternatif menarik bagi pengembang yang membutuhkan opsi beyond Raspberry Pi~\cite{orangepi2024pi5pro}.

Pemilihan Raspberry Pi 5 dan Orange Pi 5 Pro dalam studi ini didasarkan pada kemudahan ketersediaan di toko online lokal Indonesia, kemudahan dukungan teknis di komunitas, serta spesifikasi yang secara praktis memenuhi kebutuhan eksperimen ERP dan edge computing. Kedua SBC tersebut adalah produk yang paling umum ditemukan dan menawarkan kombinasi balance antara performa, harga serta kompatibilitas perangkat lunak.

Raspberry Pi 5 dan Orange Pi 5 Pro dapat dibandingkan dari sisi spesifikasi utama berikut:

\begin{table}[ht]
\centering
\begin{tabular}{|l|c|c|}
\hline
\textbf{Spesifikasi} & \textbf{Raspberry Pi 5} & \textbf{Orange Pi 5 Pro} \\
\hline
Prosesor & Cortex-A76 Quad-core 2.4GHz & RK3588S Octa-core 2.4GHz\\
RAM Maksimal & 8GB LPDDR4X & 32GB LPDDR4/LPDDR5 \\
Output Video & Dual HDMI 4K & Output 8K HDMI, DP \\
Storage & microSD, PCIe M.2 SSD & microSD, M.2 NVMe SSD \\
Port Ekspansi & PCIe, GPIO, CSI/DSI & PCIe, GPIO, M.2, DSI/CSI \\
Harga Pasaran & Rp1,200,000–1,800,000 & Rp1,800,000–2,700,000 \\
\hline
\end{tabular}
\caption{Perbandingan utama Raspberry Pi 5 dan Orange Pi 5 Pro}
\label{tab:pi5-opi5pro}
\end{table}

\section{Penelitian Terkait Development Aplikasi pada SBC}

Rangkuman dari Putra2019 dan Arief 2018 (mungkin 4-7 paragraf, jangan panjang panjang)

\section{Metode Evaluasi}

\subsection{Evaluasi Teknis dan Performa}

\begin{enumerate}
    \item Metrik Performa Aplikasi Enterprise
    \item Resource Utilization Metrics
    \item Load Testing dan Stress Testing
    \item Benchmarking Methodology
\end{enumerate}

\subsection{Evaluasi Finansial}

\begin{enumerate}
    \item Konsep Total Cost of Ownership (TCO)
    \item Komponen Biaya Infrastruktur Cloud (VPS)
    \item Komponen Biaya Infrastruktur On-Premise (SBC)
    \item Model Perhitungan dan Perbandingan TCO
\end{enumerate}

\section{Kriteria Kelayakan Development ERP}

\subsection{Kelayakan Teknis}

Konten: Minimum system requirements untuk ERP; performance thresholds (acceptable response time, minimum throughput); stability dan reliability requirements; scalability requirements; security considerations

\subsection{Kelayakan Finansial}

Konten: Cost-benefit analysis framework; acceptable ROI/payback period; affordability thresholds; comparison dengan existing solutions (SaaS ERP pricing); value for money assessment

\subsection{Trade-off Analysis Framework}

Konten: Multi-criteria decision making (MCDM) methods; Pareto analysis; weighted scoring models; decision matrix untuk mempertimbangkan performa vs cost vs other factors